\begin{proof}

\begin{enumerate}[label={$\langle1\rangle$\arabic*.},leftmargin=*]
\item We apply \cite[Chapter~VI, Theorem~1.1 and Table~10, pp.~243--244]{BHPV2004E} to the smooth compact
normalization components, using the identity $e(S_i)=K_{S_i}^2$ proved
in P4.6. The relevant non-Kähler classification is stated in
\cite[Theorem~4.2 and proof, pp.~17--21; Theorem~5.1 and proof, pp.~27--29]{OVV2020}: a minimal non-Kähler
surface is elliptic or of class VII, and the reduced fibres of a
non-Kähler elliptic surface are smooth elliptic curves.

\item For a component with nonempty conductor, a Kähler form $\omega$ would
give $K_{S_i}\cdot\omega=-D_i\cdot\omega<0$. No positive power of
the canonical bundle can then have a nonzero section. The classification
makes a Kähler component rational or ruled. If it is obtained by $k$
blowups of a ruled surface over a genus-$g$ curve, then
\[
 e=4(1-g)+k,\qquad K^2=8(1-g)-k.
\]
Equality gives $k=2(1-g)$, so $g\le1$ and the genus-one case is minimal.
A rational component has $\chi(\OO)=1$, hence $e=6$ here, including
the case with minimal model $\Pp^2$.

\item A minimal non-Kähler elliptic surface has $e=K^2=0$, by its smooth
elliptic reduced fibres. A minimal class-VII surface has $b_1=1$ and
$p_g=0$, so $\chi(\OO)=0$ and $K^2=-e$, with $e=b_2\ge0$.
Blowups preserve $\chi(\OO)$ and change $(e,K^2)$ by $(1,-1)$.
Consequently $e=K^2$ forces minimality and $e=0$ in all these
non-Kähler cases. The class-VII case has $b_2=0$. Theorem 1 of
Buonerba--Bogomolov--Kurnosov \cite[Theorem~1, p.~1]{BBK2019} supplies a finite unramified Hopf or
Inoue cover. The conductor excludes an Inoue cover: Inoue surfaces
have no compact curves, by Inoue's theorem as stated and credited in
\cite[Theorem~3.1, p.~303]{KhaledTeleman2026Journal}.
The type-III family has the explicit type-II unramified double cover
of \cite[proof of Remark~3.5, p.~305]{KhaledTeleman2026Journal}. A finite unramified quotient of a Hopf surface has universal
cover $\C^2\setminus\{0\}$ and is therefore a Hopf surface.

\item For an irreducible singular support there is one normalization
component. P4.6 gives $T=2\chi(\OO_{S_1})$, which is two in the
rational case and zero in the other cases. P4.12 proves that the
dualizing sheaf of an irreducible support is torsion. Adjunction on
the normalization therefore makes $K_{S_1}+D_1$ torsion. Since
$D_1\ne0$, every positive plurigenus of $S_1$ vanishes: multiplying
a pluricanonical section by an appropriate positive power of the
conductor section would give a section of a torsion line bundle with
a zero, whose sufficiently high power is a nonzero constant. Thus a
non-Kähler normalization in this irreducible case is of class VII,
and hence is the Hopf case already obtained.

\item For a reducible connected fibre, each normalization component meets
another global irreducible component. Pulling back the divisor of
that component gives a nonzero effective divisor on $S_i$ with zero
real first Chern class. A Kähler form would pair positively with it,
contradicting this zero class. Hence every component is non-Kähler,
has $\chi(\OO)=0$, and $T=0$. The conductor components are disjoint
and smooth, because all their nodes would come from triple points.
For any conductor component $C$, adjunction and $K+D\equiv0$ give
$2g(C)-2=(K+C)\cdot C=0$, since the other conductor components
are disjoint. Thus the double curves are smooth elliptic, and their
branch maps are the unramified maps already described in P4.6.

\item Finally let the support $G$ of a multiple fibre $mG$ be smooth. The
normalized cyclic base change $t=u^m$ is locally a disjoint union of
smooth charts: the original equation is $t=z^m$ times a holomorphic
unit, whose root can be extracted locally. These charts show that
the new central fibre maps finitely and unramifiedly to $G$, and
that the new map to the $u$-disc is a proper submersion. Its nearby
fibres are the original tori. Catanese's theorem
\cite[Theorem~2.1]{Catanese2004} states that every
deformation of a complex torus is a complex torus; its definition
uses precisely a proper holomorphic submersion over a connected
disc, with no Kähler hypothesis on the total space. We have verified
all these hypotheses. The new central fibre is therefore a torus;
Ehresmann's theorem also shows that it is connected. This gives
the required finite unramified cover of $G$.

\end{enumerate}
\end{proof}
